\usepackage{etoolbox} %programming facilities for Symbolverzeichnis
\usepackage[ngerman]{babel} %Sorachdefinition für alles
\usepackage[backend=biber,style=ieee]{biblatex} %Bibliotheksmanagement
%\usepackage[onehalfspacing]{setspace} %% Zeilenabstände, singlespacing 1-fach, onehalfspacing 1,5-fach, doublespacing 2-fach 
\usepackage[T1]{fontenc}
\usepackage{textcomp} % Verhindert Warnungen bei bestimmten Sonderzeichen
\usepackage{lmodern}
\usepackage[scaled]{helvet}   % Lädt Helvetica (phv)
\renewcommand{\familydefault}{\sfdefault} 
\usepackage[italic]{mathastext} 
\Mathastext
\usepackage{graphicx}
\graphicspath{{Bilder/}} %Bildpfad, nötig für
\usepackage{csquotes} %nötig für biblatex und \textquote (Anführungszeichen)
\usepackage{fancyhdr} % Kopf/Fußzeile in SEHR GEIL
\usepackage{extramarks}
\usepackage[a4paper,tmargin=2cm, lmargin=3.2cm,rmargin =2cm,bmargin=3cm]{geometry} %15,5 cm Platz zum schreiben
\usepackage{tabularx} %Bestes Tabellenpackage
\usepackage{icomma} %kein Platz nach einem Komma in mathmode
\usepackage{lastpage} % Ausgabe der letzten Seite
%\usepackage{stix} % Nötig für HAW: Schriftarten
\usepackage{amsmath} %Vernünftiger Mathmode bit \begin{align}
\usepackage{xfrac} % Darstellung von Equations im Text
\usepackage{tikz} % Grafikzeichenpackage
\usepackage{pdfpages} %Vernünftige Einbindung von PDFs auch als ganze Seiten mit Beibehalten des Formats
\usepackage{svg} %Bilder im svg Format 
\usepackage[table]{xcolor} %Vernünftige Farben in Tabellen
\usepackage[nottoc,notbib]{tocbibind} %Einbindung der Abbildungsverzeichnise ins TOC
\usepackage[hang,font=small]{caption} %Mehrzeilige Darstellung von Überschriften
\usepackage[hang]{footmisc} % Mehrzeiliges Einrücken von Fußnoten
\usepackage{tablefootnote} % Fußzeilen in Tabellen
\usepackage[absolute]{textpos} %Wichtig für Titelseite
\usepackage[section]{placeins} %\Floatbarrier
\usepackage[normalem]{ulem} %Texte durchstreichen mit \sout bzw unterschtreichen 
\usepackage{comment} %\begin{comment} und \end{comment} als Kommentierumgebung
\usepackage{listings} %Codedarstellung mit \begin{lstlisting}[language=python]
\usepackage[colorlinks=true,linkcolor=black,citecolor=black,urlcolor=black]{hyperref} %Hyperref als beste Refenzierungspackage


%% Text und Schriftgrößeneinstellungen 
\setlength{\parindent}{0em}% kein Einrücken
\newcommand{\changefont}{\fontsize{7}{12}\selectfont}%Kopfzeile
\newcommand{\changefontbig}{\fontsize{10}{12}\selectfont}%Kopfzeile
\setcounter{biburllcpenalty}{9000}% Kleinbuchstaben
\setcounter{biburlucpenalty}{9000}% Großbuchstaben
\widowpenalty = 10000 %Vermeidung von Satzende auf nöchster Seite
\clubpenalty = 10000 %Vermeidung von Absatzende auf nöchster Seite


%Nummerierungen innerhalb der Sections
\numberwithin{equation}{section}
\numberwithin{table}{section}
\numberwithin{figure}{section}

%%Lists in academic style
\renewcommand{\labelenumii}{\arabic{enumi}.\arabic{enumii}}
\renewcommand{\labelenumiii}{\arabic{enumi}.\arabic{enumii}.\arabic{enumiii}}
\renewcommand{\labelenumiv}{\arabic{enumi}.\arabic{enumii}.\arabic{enumiii}.\arabic{enumiv}}
\renewcommand{\labelitemi}{-}

%% KOPFZEILE DARSTELLUNG Leftmark ohne sectionnummer

\renewcommand{\sectionmark}[1]{\markboth{#1}{}}%
\renewcommand{\subsectionmark}[1]{\markright{#1}}%

% Neudefintionen von autoref Verweisen
%autoref mit Equation Darstellung Gleichung (x)
\makeatletter
  \AtBeginDocument{%
    \let\plain@equationautorefname\equationautorefname
    \def\equationautorefname{\plain@equationautorefname\@autoref@insert@tagform}%
    \def\@autoref@insert@tagform~#1\null{~(#1\null)}%
  }
\makeatother

% Kapitelverweise mit Abschnitt statt Standard
\addto\extrasngerman{\def
\subsectionautorefname{Unterkapitel}
}

\addto\extrasngerman{\def
\subsubsectionautorefname{Unterkapitel}
}

\addto\extrasngerman{\def
\sectionautorefname{Kapitel}
}


% Ändert den Text, den \autoref verwendet

\renewcommand{\lstlistingname}{Code-Ausschnitt}
\def\lstnumberautorefname{Zeile}

% --- Farben für das Spyder-Theme definieren ---
% Basierend auf deinem Screenshot
\definecolor{spyderBlue}{HTML}{0000FF}    % Keywords (import, from)
\definecolor{spyderGreen}{HTML}{008000}   % Strings und Docstrings
\definecolor{spyderGrey}{HTML}{999999}    % Kommentare
\definecolor{spyderCyan}{HTML}{9a0aab}    % Klassen, spezielle Funktionen
\definecolor{spyderNumber}{HTML}{B8860B}  % Ein gut lesbarer Gelb/Gold-Ton
\definecolor{spyderself}{HTML}{b80b1f}

% --- Den Listing-Stil definieren ---
\lstdefinestyle{spyderstyle}{
    language=Python,
    backgroundcolor=\color{white},   
    commentstyle=\color{spyderGrey},
    keywordstyle=\color{spyderCyan},
    stringstyle=\color{spyderGreen},
    numberstyle=\tiny\color{darkgray}, 
    basicstyle=\ttfamily\footnotesize\color{black},
    emphstyle=\color{spyderCyan},
    emph={[2]int, str, float, list, dict, set, tuple, print, range, open, len},
    emphstyle={[2]\color{spyderBlue}\bfseries},
    emph={[3]self},
    emphstyle={[3]\color{spyderself}\bfseries},
    escapeinside={(*@}{@*)},
    literate=
      *{0}{{{\color{spyderNumber}0}}}1
       {1}{{{\color{spyderNumber}1}}}1
       {2}{{{\color{spyderNumber}2}}}1
       {3}{{{\color{spyderNumber}3}}}1
       {4}{{{\color{spyderNumber}4}}}1
       {5}{{{\color{spyderNumber}5}}}1
       {6}{{{\color{spyderNumber}6}}}1
       {7}{{{\color{spyderNumber}7}}}1
       {8}{{{\color{spyderNumber}8}}}1
       {9}{{{\color{spyderNumber}9}}}1
       {.}{{{\color{spyderNumber}.}}}1
       {e-}{{{\color{spyderNumber}e-}}}2
       {e+}{{{\color{spyderNumber}e+}}}2
       {E-}{{{\color{spyderNumber}E-}}}2
       {E+}{{{\color{spyderNumber}E+}}}2,,
    breakatwhitespace=false,         
    breaklines=true,                 
    captionpos=b,                    
    keepspaces=true,                 
    numbers=left,                    
    numbersep=5pt,                  
    showspaces=false,                
    showstringspaces=false,
    showtabs=false,                  
    tabsize=4}

% Den Stil als Standard für alle Listings setzen
\lstset{style=spyderstyle}


%% Anhangsverzeichnis

%Aufsplittung Inhalts und Anhangsverzeichnis 
\makeatletter% 
\newcommand*{\maintoc}{% Hauptinhaltsverzeichnis
  \begingroup
    \@fileswfalse% kein neues Verzeichnis öffnen
    \renewcommand*{\appendixattoc}{% Trennanweisung im Inhaltsverzeichnis
      \value{tocdepth}=-10000 % lokal tocdepth auf sehr kleinen Wert setzen
    }%
    \tableofcontents% Verzeichnis ausgeben
  \endgroup
}
\newcommand*{\appendixtoc}{% Anhangsinhaltsverzeichnis
  \begingroup
    \edef\@alltocdepth{\the\value{tocdepth}}% tocdepth merken
    \setcounter{tocdepth}{-10000}% Keine Verzeichniseinträge
    \renewcommand*{\contentsname}{% Verzeichnisname ändern
      Anhangsverzeichnis}%
    \renewcommand*{\appendixattoc}{% Trennanweisung im Inhaltsverzeichnis
      \setcounter{tocdepth}{\@alltocdepth}% tocdepth wiederherstellen
    }%
    \tableofcontents% Verzeichnis ausgeben
    \setcounter{tocdepth}{\@alltocdepth}% tocdepth wiederherstellen
  \endgroup
}
\newcommand*{\appendixattoc}{% Trennanweisung im Inhaltsverzeichnis
}
\g@addto@macro\appendix{% \appendix erweitern
  \clearpage% Neue Seite
  \addcontentsline{toc}{section}{\appendixname}% Eintrag ins Hauptverzeichnis
  \addtocontents{toc}{\protect\appendixattoc}% Trennanweisung in die toc-Datei
}
\makeatother



%% Kurzbefehle für eigene Designcommands

%Römische Zahlen
\newcommand{\uproman}[1]{\uppercase\expandafter{\romannumeral#1}}
\newcommand{\lowroman}[1]{\romannumeral#1\relax}


%Kreis um z.B. Zahl
\newcommand\circlearound[1]{\unitlength1ex\begin{picture}(2.1,2.1)%
\put(0.5,0.5){\ \circle{2.1}}\put(0.5,0.5){\ \makebox(0,0){#1}}\end{picture}}


%%plus und minus in kleiner
\newcommand{\plus}{\raisebox{.4\height}{\scalebox{.6}{+}}}
\newcommand{\minus}{\raisebox{.4\height}{\scalebox{.8}{-}}}


%% Kurzbefehle fürs Einfügen von Bildern und Svgs

\newcommand*{\newfigure}[4][]{%
    \begin{figure}[ht!]%
        \centering%
        \includegraphics[width=#2\textwidth]{#3}%
        \caption[#1]{#4}%
         \label{fig:#3}%
        \ifx&#1&%
            % No short caption provided, use full caption for list of figures
            \addcontentsline{lof}{figure}{\protect\numberline{\thefigure}#4}%
        \fi%
    \end{figure}%
    \FloatBarrier
    }


\newcommand*{\newsvg}[4][]{%
    \begin{figure}[ht!]%
        \centering%
        \includesvg[width=#2\textwidth]{#3}%
        \caption[#1]{#4}%
        \label{fig:#3}%
        \ifx&#1&%
            % No short caption provided, use full caption for list of figures
            \addcontentsline{lof}{figure}{\protect\numberline{\thefigure}#4}%
        \fi%
    \end{figure}%
        \FloatBarrier
}

\newcommand*{\newmultifig}[4]{%
\begin{figure}[ht!]
\centering
\begin{minipage}[t]{0.46\textwidth}
    \centering
    \includegraphics[width=\textwidth]{#1}
    \caption{#2}
    \label{fig:#1}   
\end{minipage}
\hspace{0.03\textwidth}
\begin{minipage}[t]{0.46\textwidth}
    \centering
    \includegraphics[width=\textwidth]{#3}
    \caption{#4}
    \label{fig:#3}
    \end{minipage}
\end{figure}
    \FloatBarrier
}


\newcommand{\newparagraph}[1]{
\paragraph{#1}\label{#1}\mbox{}\\[1ex]
}

